\documentclass{article}
\usepackage{amsmath,amssymb}
\usepackage{xurl}
\usepackage[hidelinks]{hyperref}
% Shared commands for notes in docs/paper-gaps/.

\DeclareUnicodeCharacter{2080}{\ifmmode{}_{0}\else\textsubscript{0}\fi}
\DeclareUnicodeCharacter{2081}{\ifmmode{}_{1}\else\textsubscript{1}\fi}
\DeclareUnicodeCharacter{2082}{\ifmmode{}_{2}\else\textsubscript{2}\fi}
\DeclareUnicodeCharacter{2099}{\ifmmode{}_{n}\else\textsubscript{n}\fi}

\newcommand{\Lean}{\textsc{Lean}}
\ExplSyntaxOn
\NewDocumentCommand{\leanid}{m}{
  \ifmmode
    \textsf{\detokenize{#1}}
  \else
    \group_begin:
    \urlstyle{sf}
    \tl_set:Nn \l_tmpa_tl {#1}
    \tl_replace_all:Nnn \l_tmpa_tl {\_} {_}
    \exp_args:NV \path \l_tmpa_tl
    \group_end:
  \fi
}
\ExplSyntaxOff
\providecommand{\tr}{\operatorname{tr}}
\newcommand{\tnleanIssueBase}{https://github.com/LionSR/TNLean/issues/}
\newcommand{\tnleanPrBase}{https://github.com/LionSR/TNLean/pull/}
\newcommand{\tnleanDocBase}{https://sirui-lu.com/TNLean/docs/}
\newcommand{\ghissue}[1]{\href{\tnleanIssueBase#1}{GitHub issue \##1}}
\newcommand{\ghpr}[1]{\href{\tnleanPrBase#1}{PR \##1}}
\newcommand{\doclink}[1]{\href{\tnleanDocBase#1.html}{\path{#1}}}

% Dirac notation and the identity operator, as in blueprint/src/macros/common.tex.
% \providecommand keeps note-local definitions authoritative.
\usepackage{dsfont}
\providecommand{\ket}[1]{|#1\rangle}
\providecommand{\bra}[1]{\langle #1|}
\providecommand{\braket}[2]{\langle #1 | #2 \rangle}
\providecommand{\ketbra}[2]{|#1\rangle\!\langle #2|}
\providecommand{\Id}{\mathds{1}}
\providecommand{\one}{\mathds{1}}
\providecommand{\C}{\mathbb{C}}
\providecommand{\MN}[1]{M_{#1}(\C)}
\providecommand{\MD}{M_D(\C)}

% Verdict marker: \gapnote{<kind>}{<status>}. Kind is the policy
% classification (clarification, local-correction, scope-restriction,
% unfaithful, false-source, open-gap); status is open, wip (elimination
% underway), resolved, or historical. Machine-read by the site index;
% prints a verdict line.
\newcommand{\gapnote}[2]{\par\noindent\textbf{Verdict:} #1 (#2).\par}

\providecommand{\leanid}[1]{\textsf{\detokenize{#1}}}
\providecommand{\gapnote}[2]{\par\noindent\textbf{Verdict:} #1 (#2).\par}
\title{Principal Density Decay at the Rounded Scanner Scales}
\author{}
\date{2026-10-14}
\begin{document}
\maketitle
\gapnote{scope-restriction}{open}

\section{Numerical conclusion}

Fix scanner exponents, a nonnegative coefficient $C$, a real logarithmic
power $C_\ell$, a weight exponent $\omega$ and a target exponent $r$.
Assume the two strict inequalities
\begin{align}
    e-\mu+\nu&<-r,&
    \frac{\omega-(1-\ell-\mu)}4+\nu&<-r.
    \label{eq:density_scope_margins}
\end{align}
These are additional requirements beyond the ordinary scanner exponent
conditions. There is a threshold $N\geq2$, depending on these fixed
quantities, such that every integer $n\geq N$ and every
$1\leq W\leq n^\omega$ satisfy $K\geq1$ and
\begin{align}
    \frac{C}{n^{-\nu}}
      \left(\frac{n^e}{m}+a^{1/4}(\log n)^{C_\ell}\right)
    &\leq n^{-r},
    \label{eq:density_scope_bound}
\end{align}
where $L=\lfloor n^{1-\ell}\rfloor$, $m=\lfloor n^\mu\rfloor$,
$K=L\mathbin{\mathrm{div}}(8m)$ and $a=W/K$.
Both floors and the natural-number division are retained. The constant
and logarithmic power precede the threshold; the threshold precedes
the scale, weight and any physical instance to which it is applied.
No sign condition on $C_\ell$ is required.
\footnote{The numerical declarations are
  \leanid{TNLean.PEPS.AreaLaw.Scan.ScannerExponents.eventually_principalDensity_le_rpow}
  and \leanid{TNLean.PEPS.AreaLaw.Scan.ScannerExponents.exists_principalDensity_le_rpow}.}

The proof uses only the rounded-scale estimates
$n^e/m\leq C_mn^{e-\mu}$ and
$n^{1-\ell-\mu}\leq C_KK$, followed by the fact that every fixed
logarithmic power is smaller than a positive power of $n$.
No terminal energy, entropy, density or selected-point conclusion is
assumed. The separation ceiling $D=\lceil n^\kappa\rceil$ does not occur
in (\ref{eq:density_scope_bound}).

\section{Final source parameters}

At $e=2/1000000$, $\mu=1/4$, $\nu=1/1000$,
$\ell=1/100000$, $\omega=1/20000$ and $\kappa=1/1000000$,
the two margins in (\ref{eq:density_scope_margins}) at $r=\nu$ are
exactly $123999/500000$ and $37097/200000$. Thus the conclusion
(\ref{eq:density_scope_bound}) holds with $r=\nu$.
\footnote{The declarations are
  \leanid{TNLean.PEPS.AreaLaw.Scan.final_principalDensity_margins}
  and \leanid{TNLean.PEPS.AreaLaw.Scan.exists_final_principalDensity_le_rpow}.
  Source: \cite{OpenAI2026AreaLaw}, \path{08-scanner.tex},
  \texttt{scanner:scales}, lines 32--41;
  \texttt{scanner:density-output}, lines 380--384;
  \texttt{scanner:final-parameters}, lines 839--865;
  revision \texttt{adc7f1241b42e322a6451854ab7e4b4c146bf78a}.}

\section{Remaining physical obligations}

This is the numerical estimate for the principal density expression.
It does not establish that an actual selected entropy defect is bounded
by that expression. That implication still requires the selected-charge
entropy telescope, the actual terminal norm comparisons, the
original-state target and shell entropy bounds and the treatment of any
additional errors, all for the same physical data and label sequence.
It does not construct those inputs, prove the remaining transport
smallness or geometric inequalities, estimate the other bootstrap
errors, or complete the mutual-information limit and continuity transfer.
Consequently it does not assert Proposition~9.6, a full physical scanner,
or the area law. The scope restriction remains open.

\bibliographystyle{alpha}
\bibliography{references}
\end{document}
